%
\tikzsetnextfilename{fig_small_gain_theorem}
\centering
%, line width=0.7
\tikzstyle{comp} = [circle, draw, minimum size = 0.4cm] % comparator style
\tikzstyle{block} = [rectangle, draw, minimum width = 0.6cm, minimum height = 0.6cm]  % bloc style

\newcommand{\comp}[3][]{ % comparator command
	\node[comp] (#2) at #3 {};
	\draw (#2) -- + (45:0.2) -- + (45:-0.2);
	\draw (#2) -- + (-45:0.2) -- + (-45:-0.2);
	\ifthenelse{
		\isempty{#1}{}
	}{}
	{
		\ifthenelse{
			\equal{#1}{a}
		}{
			\node[above right] at (#2) {$-$\phantom{p}};
		}{}
		\ifthenelse{
			\equal{#1}{b}
		}{
			\node[below right] at (#2) {$-$\phantom{l}};
		}{}
	}
}
\begin{tikzpicture} [node distance = 1cm]
	
	\node[block] (M) {$M$};
	\node[block, below of = Z] (D) {$\Delta$};
	
	\comp{c1}{($(M) + (-1.5, 0)$)};
	\comp{c2}{($(D) + (1.5, 0)$)};
	
	\draw[-{Latex[length=3mm]}] (c1.east) -- (M.west);
	\draw[-{Latex[length=3mm]}] (M.east) -| (c2.north);
	\draw[-{Latex[length=3mm]}] (c2.west) -- (D.east);
	\draw[-{Latex[length=3mm]}] (D.west) -| (c1.south);
	
	\draw[{Latex[length=3mm]}-] (c1.west) --+ (-.75,0) node[near end, above] {$w_1$};
	\draw[{Latex[length=3mm]}-] (c2.east) --+ (.75,0) node[near end, above] {$w_2$};
	
\end{tikzpicture}
\captionof{figure}[Schéma bloc du théorème du petit gain \parencite{zhou_robust_1995}]{Théorème du petit gain \parencite{zhou_robust_1995}.}
\label{fig_small_gain_theorem}
